\subsubsection{高水準設計}

\term{高水準設計}{High-Level Design} は、システムの主要構成要素が互いにどう関わり、外部環境とどうやり取りするかを決める設計である。外部環境には、利用者、外部システム、デバイス、ネットワーク、データベース、運用環境などが含まれる。

高水準設計では、主に次を扱う。
\begin{itemize}
  \item システムが反応すべき外部イベントやメッセージ。
  \item システムが生成すべきイベントやメッセージ。
  \item イベントやメッセージのデータ形式とプロトコル。
  \item 入力と出力の順序、タイミング、依存関係。
  \item 端から端までの取引やイベント列の流れ。
  \item データをどのように保存し、管理するか。
\end{itemize}

高水準設計は、アーキテクチャがある場合、その枠内で行われる。たとえば、アーキテクチャが「非同期メッセージでやり取りする」と決めていれば、高水準設計ではその方針に従って、具体的なメッセージ種類、順序、例外時の扱いを決める。
